% Template BÔNUS — estilo "empilhado" (cabeçalho à esquerda, com Resumo).
% Mesmo princípio ATS-friendly: 1 coluna, sem foto, texto linear.
\documentclass[a4paper,10pt]{article}
\usepackage[utf8]{inputenc}
\usepackage[T1]{fontenc}
\usepackage[portuguese]{babel}
\usepackage{geometry}
\usepackage{parskip}
\usepackage{microtype}
\usepackage{enumitem}
\usepackage{titlesec}
\usepackage{hyperref}
\usepackage{xurl}
\geometry{top=1.4cm, bottom=1.4cm, left=1.6cm, right=1.6cm}
\setcounter{secnumdepth}{0}
\setlist[itemize]{leftmargin=1.1em, itemsep=0.15em, topsep=0.2em, parsep=0em, partopsep=0em}
\pagestyle{empty}
\hypersetup{colorlinks=true, linkcolor=black, urlcolor=black}

% Títulos de seção em negrito, com linha embaixo (estilo "empilhado")
\titleformat{\section}{\large\bfseries}{}{0em}{}[\titlerule\vspace{0.3ex}]
\titlespacing*{\section}{0pt}{0.9em}{0.35em}

% Entrada: cargo/título em negrito + subtítulo em itálico (empresa · local · datas)
\newcommand{\entry}[2]{\noindent\textbf{#1}\\ \textit{#2}\par\vspace{0.15em}}

\begin{document}

% --- CABEÇALHO (alinhado à esquerda) ---
{\Huge\textbf{Seu Nome}}\\[3pt]
{\large Seu Cargo ou Área (ex.: Estudante de Gestão em Saúde)}\\[3pt]
Sua Cidade, UF \textbullet{} Disponível para remoto\\
+55 51 90000-0000 \textbullet{} \href{mailto:seu.email@exemplo.com}{seu.email@exemplo.com}\\
\href{https://www.linkedin.com/in/seu-usuario}{LinkedIn} \textbullet{} \href{https://github.com/seu-usuario}{GitHub} (opcional)

\section{RESUMO}
Um parágrafo curto (3 a 4 linhas) com seu perfil, foco e objetivo. Use palavras-chave reais
da área/vaga — sem frases genéricas. Ex.: "Estudante de Gestão em Saúde com experiência em
indicadores e processos administrativos, buscando estágio na área de qualidade hospitalar."

\section{EDUCAÇÃO}
\entry{Seu Curso (em andamento)}{Nome da Universidade \textbullet{} Cidade, UF \textbullet{} Ano -- Conclusão prevista}
\entry{Ensino Médio Completo}{Nome da Escola \textbullet{} Cidade, UF \textbullet{} Ano}

\section{EXPERIÊNCIA}
\entry{Seu Cargo}{Nome da Empresa \textbullet{} Cidade, UF \textbullet{} Mês Ano -- Mês Ano}
\begin{itemize}
    \item Comece com um substantivo de ação e descreva o que você fez.
    \item Seja específico; quantifique o resultado quando possível.
    \item No máximo 3 bullets por entrada.
\end{itemize}

\section{PROJETOS E VOLUNTARIADOS}
\entry{Nome do Projeto / Liga / Iniciação Científica}{Instituição \textbullet{} Cidade, UF \textbullet{} Mês Ano -- Mês Ano}
\begin{itemize}
    \item Atividades acadêmicas, extensão, monitoria, liga, voluntariado.
\end{itemize}

\section{HABILIDADES}
\textbf{Ferramentas:} Ex.: Excel, Power BI, Google Workspace\\
\textbf{Conhecimentos:} Ex.: indicadores em saúde, gestão de processos\\
\textbf{Certificações:} Ex.: curso livre relevante (ano)\\
\textbf{Idiomas:} Português (Nativo), Inglês (Intermediário)

\end{document}
